% Study 2: supplementary methods, the relation to the CogSci 2025 analysis, and the results that
% no other part of the paper prints. Numbers come only from macros in generated/numbers.tex.
\label{sec:suppB-study2}

Study~2 asked whether the AI label on an original reaches ratings of a second work that Alex was said to have made from it with his design and sketching tools. It was the first of the five studies to be run, from 23 November to 5 December 2024, and our CogSci 2025 paper,\cite{kamper2025originalsin} a peer-reviewed proceedings paper, reported it as its Study~1. A registration is a dated record of a study's hypotheses and analysis plan, unrelated to copyright registration. Study~2's was filed on the Open Science Framework (OSF) on 8 October 2026, after its data had been collected and analyzed. Its registered verdicts are in Supplementary Table~\ref{tab:supp-verdicts}, the label's effect averaged over levels is in Supplementary Table~\ref{tab:supp-labeleffects}, and the main text reports both. Supplementary Note~6 gives its registered test of dilution (H3), and Supplementary Note~7 the analyses of its explanations pooled with those of Studies 4 and 5. This section gives the materials and procedure, the sample, how the registered analysis differs from the earlier one, and the results that no other part of the paper prints.

\subsection*{Materials}

Every participant saw the same ten original images. All ten were generated with OpenAI's DALL·E, an AI image generator, so the Human label described below was a cover story (Section~\ref{sec:suppA-registrations} gives the evidence for the generator). Five originals show a human or humanoid figure and five do not, and the registered model includes this image kind as a factor. Each original came with three second works: None, the original's own image file shown again unchanged; Slight, described in the registration as adjustments of color and contrast; and Dramatic, substantial changes. We made the Slight and Dramatic versions by hand in image-editing software, and no AI tool changed any image (Section~\ref{sec:suppA-registrations}). Of the ten Slight versions, \XSImageCheckSlightReframedN{} were also cropped or resized (Supplementary Note~4), although the CogSci 2025 paper describes the Slight edits only as adjustments of hue, coloration and contrast. Studies 2, 4 and 5 showed the same 30 images.

Participants read the label once, at the start. The Human label read: `There is a digital artist named Alex. Alex is a skilled artist who creates digital artworks. Alex only uses graphic design software and digital sketching tools.' The AI label read: `There is a digital artist named Alex. Alex is a skilled artist who creates digital artworks. Alex develops and inputs prompts into an artificial intelligence (AI) system. The AI system generates digital artworks based on those prompts.'

Before each original, the Human group read `Alex had a new idea for a digital artwork. Alex created a new piece of digital artwork only using graphic design software and digital sketching tools.' The AI group read `Alex had a new idea for a work of digital artwork. Alex prompted the artificial intelligence (AI) system for a new piece of digital artwork, which created a new digital artwork.' Both went on: `Click the arrow to see the digital artwork created. Note that you will be required to stay on the page for at least 8 seconds before proceeding.'

Both groups read the same description of each second work. Before the first of the three, it was `Alex looked at the first artwork, and decided to produced a secondary digital artwork. Using the first work as a base, he produced the second digital artwork with his graphic design software and digital sketching tools.' Before each of the other two, it was `Later, Alex looked at the very first artwork again. Alex decided to produce a different secondary digital artwork. Using the very first artwork as a base, he produced a different secondary digital artwork with his graphic design software and digital sketching tools.' Each ended `Click the arrow to see the digital artworks:'. These quotations reproduce the survey exactly, slips included. In both groups, then, Alex made every second work himself; only the account of the original differed. No description said how much had been changed or how long anything took.

Three features of this wording belong to Study~2 alone. Its Human label says that Alex `only uses graphic design software', where Studies 1, 4 and 5 say that he `uses his graphic design skills' and Study~3 names the tools. Its AI vignette before each original says nothing about Alex revising his prompts, unlike the other four studies. And its second-work descriptions name Alex's `graphic design software', where Studies 4 and 5 name `his own graphic design skills'.

On the screen after each pair of original and second work, which showed neither image, participants gave three ratings on five-point scales, always in the order below, each scored 1 to 5 by the position of the answer chosen. The transformativeness item read `How transformative is the second work compared to the first work? Transformativeness is the degree the second work adds new expression, meaning, or message to the first work.' Its answers ran from `Not at all transformative' through slightly, moderately and very transformative to `Highly transformative'. The essence item read `How much has the essence of the first work changed in the second work?', with answers from `No change in essence' through very little, moderate and significant change to `Complete change in essence'. The creativity item read `How much creativity was involved in the production of the second work?', with answers from `No creativity involved' through minimal, moderate and significant creativity to `A great deal of creativity involved'. Below the three ratings, a required text box for each asked `Please explain your reasoning for your rating on transformativeness for your responses:' (or on essence, or on creativity). Each explanation was therefore written after its rating had been given, although a rating could still be changed before the participant left the screen. A script blocked copying and pasting in 21 of the 90 explanation boxes and not in the others.

Study~2 is the only study whose essence and creativity items carried no definition. Studies 4 and 5 defined essence as `the core sense and feeling that one has while observing a work of art that makes that work unique', and the other four studies defined creativity as `the presence of at least a minimal original, intellectual contribution to the creation of the work'. The CogSci 2025 paper quotes all three items with definitions, but the Study~2 survey file shows a definition on the transformativeness item only. Comparisons of creativity ratings between Study~2 and the other studies are therefore exploratory, as Study~2's registration states. Two transformativeness items, both for one original (at None and at Slight), had six answer points, a reworded question and no definition. As registered, their \STwoNSixPointDropped{} ratings and explanations were left out. One creativity item asked about `creating the second work', and three essence items had slips in their answer labels; scoring by answer position is unaffected.

\subsection*{Procedure}

Participants took part online in a Qualtrics survey. By our account, participants were students in the UCLA Psychology subject pool who took part for course credit. Every analyzed record carries the code that the survey sent, at its end, to the pool's page for granting credit. The consent form, the same in all five studies, also offered hourly pay to people who signed up outside the pool, and no record shows that anyone took that route. The survey assigned each participant at random to one label for the whole session, presenting each label equally often. It also set a quota of 125 participants for each label group, a planned number that neither stopped anyone nor moved anyone to the other label. Each participant saw all ten originals in random order. For each, the vignette came first, then the original alone for at least 8 seconds, then the three second works. Each second work was introduced by its description and then shown beside the original, under the headings `FIRST WORK' and `SECOND WORK', for at least 10 seconds; its ratings and explanations followed on the next screen. The first of the three levels was drawn at random and the other two followed in random order; the order was not recorded. Each participant thus rated 30 second works, apart from \STwoNLackingOneTrial{} whose records lack the Slight trial of one original.

A question about the task (`What is your task in this experiment?') came before the label and allowed one retry, so it ended no session. An arithmetic question (`What does 8+7 equal?') appeared at a varying point among the originals. After the last rating, a closing screen asked first whether the participant had completed the experiment seriously, without responding randomly, then whether any part of the procedure had been unclear, and then about race, gender, age, native language and formal artistic training. No question asked participants to recall the label, a check known as a manipulation check. The debriefing was the same as in the other four studies; Section~\ref{sec:suppA-study1} quotes it and corrects its statement of the law. Participants read it only after their last rating.

\subsection*{Exclusions and sample}

The registered exclusion rules, applied in order, left \STwoN{} of the \STwoNRecords{} records in the survey export. There were no survey test runs. Then \STwoNUnfinishedRemoved{} unfinished sessions were removed, and \STwoNNotSeriousRemoved{} participants who gave any answer but `Yes, I did' to the seriousness question, a blank included. No one gave a wrong answer to the arithmetic question, and every record had a label. Of the \STwoN{} participants analyzed, \STwoNHuman{} saw the Human label and \STwoNAI{} the AI label. They gave \STwoNRatings{} analyzed ratings, each with its explanation. The same unfinished and non-serious records were removed in the CogSci 2025 paper, which analyzed the same participants.

The participants' mean age was \STwoAgeMean{} years (standard deviation \STwoAgeSD; range \STwoAgeMin{} to \STwoAgeMax). Of the \STwoN, \STwoNWomen{} described their gender as female, \STwoNMen{} as male, and \STwoNGenderOther{} chose `Other' or `Prefer not to say' or gave no answer. The answers on race, native language and formal artistic training are not summarized here. The CogSci 2025 paper gives age and gender figures for this study, but their counts add up to the number of finished sessions, not to the analyzed sample, so they are not reused.

\subsection*{How the registered analysis differs from the CogSci 2025 analysis}

The CogSci 2025 paper analyzed these data with mixed analyses of variance (ANOVAs). A mixed ANOVA compares condition means computed from each participant's average rating in each condition; here the label varied between participants and the level within them. Pairwise $t$~tests, comparisons of two means, followed between levels, and the two labels were compared at each level. The registered analysis differs in four ways. First, it fits a linear mixed model to every rating rather than to each participant's averages. This is a regression that gives each participant and each original its own baseline, so that differences between originals are not mistaken for differences between conditions, and it includes image kind as a factor. Second, it fixes twelve confirmatory tests and corrects them together with Holm's method, which raises each $p$-value so that the chance of calling any of the twelve significant by mistake stays at .05. A $p$-value is the chance of a difference at least this large if there were no effect. Third, it asks about the label with an equivalence test: whether the label's effect, averaged over levels, is shown to lie inside $\pm$\SourceBound{} points, half of Study~1's creativity penalty on originals, rather than whether it differs from zero. Fourth, it adds a direct test of dilution (H3; Supplementary Note~6), the account on which the label's effect is largest at None and shrinks as the edit grows, and an ordinal check.

The two analyses agree that ratings rose steeply with the level of modification (Supplementary Table~\ref{tab:suppB-s2}b). On the label, the CogSci 2025 paper found a small creativity difference averaged over levels and none at any single level. In the registered model, the ordinary two-sided test, which asks whether that averaged creativity difference is zero, is also significant ($p$~\STwoCreatPenaltyP), but it is description only. The registered verdict rests on the equivalence test instead, and it reads `not shown to be small': the test fails after Holm's correction, although its uncorrected interval lies inside the bound (Supplementary Table~\ref{tab:supp-labeleffects}). Because the averaged creativity difference had been reported before registration, this registered test re-analyzes an effect already known. The registered model's comparisons at each level are exploratory and uncorrected (Supplementary Table~\ref{tab:suppB-s2}a). Two of the three creativity comparisons fall below .05, where the earlier analysis, after its own correction, found none. They are not evidence of a penalty at any particular level.

The CogSci 2025 paper also calls this experiment pre-registered. It was not registered with any registry before data collection. Its only planning document, a template begun in November 2024 and never submitted, was last edited after collection had ended (Study~2's registration, question~1).

\subsection*{Results not reported elsewhere}

Each result below carries its status. A registered confirmatory test is one that the registration names as confirmatory, judged by its rule after Holm's correction; a result is called significant when its $p$-value is below .05. A registered secondary analysis is one the registration names outside that set. An exploratory analysis was named as exploratory or not named at all, and carries no verdict. A descriptive result, such as a mean, carries no test.

\paragraph*{Registered tests and predicted means.}
Supplementary Table~\ref{tab:suppB-s2} gives the model's predicted mean for each label at each level, that is, the mean rating the model estimates for that condition with the other factors averaged (descriptive); the label difference at each level (exploratory); and each registered test of H1 with its ordinal check. The ordinal check is a cumulative-link mixed model, which treats the five answers as ordered categories rather than evenly spaced numbers. On transformativeness and essence change it reached the same verdict as the linear model on every test it repeated. On creativity its fit did not converge, meaning that the fitting procedure stopped without settling on a solution, so the check is not completed and carries no verdict. All three registered linear models converged, and none was singular, that is, none estimated a source of variation at zero. Supplementary Note~6 gives the registered test of dilution (H3) with its ordinal check, and Supplementary Table~\ref{tab:supp-labeleffects} the test of H2.

% Study 2: predicted means by label and level, the label difference at each level, and the
% registered tests of H1 with the ordinal check.
% Every number is a macro from generated/numbers.tex; supplement/tables/README_suppB.md lists
% the pipeline tables behind them.
\begin{table}[htbp]
\centering
\footnotesize
\caption{\textbf{Study~2: predicted mean ratings, and the registered tests of modification (H1).} Ratings are on 1-to-5 scales. \textbf{a},~The linear mixed model's predicted mean for each label at each level (descriptive), and the Human-label mean minus the AI-label mean with its 95\% confidence interval (CI) and two-sided test. Study~2 registered no test at each level, so these differences are exploratory and uncorrected. \textbf{b},~The registered confirmatory tests of H1, corrected together with H2 and H3 by Holm's method across twelve tests, beside the registered ordinal check. Supplementary Note~6 gives the test of H3 and Supplementary Table~\ref{tab:supp-labeleffects} the test of H2.}\label{tab:suppB-s2}
\setlength{\tabcolsep}{4pt}
\begin{tabular}{@{}lccccc@{}}
\multicolumn{6}{@{}l}{\textbf{a}\enspace Predicted means and the label difference}\\
\toprule
Level & Human label & AI label & Human minus AI & 95\% CI & Two-sided test\\
\midrule
\multicolumn{6}{@{}l}{\textit{Transformativeness}}\\
None & \STwoTransMeanHumanNone & \STwoTransMeanAINone & \STwoTransPenaltyNoneEst & \STwoTransPenaltyNoneLo{} to \STwoTransPenaltyNoneHi & $p$~\STwoTransPenaltyNoneP\\
Slight & \STwoTransMeanHumanSlight & \STwoTransMeanAISlight & \STwoTransPenaltySlightEst & \STwoTransPenaltySlightLo{} to \STwoTransPenaltySlightHi & $p$~\STwoTransPenaltySlightP\\
Dramatic & \STwoTransMeanHumanDramatic & \STwoTransMeanAIDramatic & \STwoTransPenaltyDramaticEst & \STwoTransPenaltyDramaticLo{} to \STwoTransPenaltyDramaticHi & $p$~\STwoTransPenaltyDramaticP\\
Averaged over levels\textsuperscript{a} & \STwoTransMeanHuman & \STwoTransMeanAI & \STwoTransPenaltyEst & \STwoTransPenaltyLo{} to \STwoTransPenaltyHi & $p$~\STwoTransPenaltyP\\
\addlinespace[3pt]
\multicolumn{6}{@{}l}{\textit{Essence change}}\\
None & \STwoEssMeanHumanNone & \STwoEssMeanAINone & \STwoEssPenaltyNoneEst & \STwoEssPenaltyNoneLo{} to \STwoEssPenaltyNoneHi & $p$~\STwoEssPenaltyNoneP\\
Slight & \STwoEssMeanHumanSlight & \STwoEssMeanAISlight & \STwoEssPenaltySlightEst & \STwoEssPenaltySlightLo{} to \STwoEssPenaltySlightHi & $p$~\STwoEssPenaltySlightP\\
Dramatic & \STwoEssMeanHumanDramatic & \STwoEssMeanAIDramatic & \STwoEssPenaltyDramaticEst & \STwoEssPenaltyDramaticLo{} to \STwoEssPenaltyDramaticHi & $p$~\STwoEssPenaltyDramaticP\\
Averaged over levels\textsuperscript{a} & \STwoEssMeanHuman & \STwoEssMeanAI & \STwoEssPenaltyEst & \STwoEssPenaltyLo{} to \STwoEssPenaltyHi & $p$~\STwoEssPenaltyP\\
\addlinespace[3pt]
\multicolumn{6}{@{}l}{\textit{Creativity}}\\
None & \STwoCreatMeanHumanNone & \STwoCreatMeanAINone & \STwoCreatPenaltyNoneEst & \STwoCreatPenaltyNoneLo{} to \STwoCreatPenaltyNoneHi & $p$~\STwoCreatPenaltyNoneP\\
Slight & \STwoCreatMeanHumanSlight & \STwoCreatMeanAISlight & \STwoCreatPenaltySlightEst & \STwoCreatPenaltySlightLo{} to \STwoCreatPenaltySlightHi & $p$~\STwoCreatPenaltySlightP\\
Dramatic & \STwoCreatMeanHumanDramatic & \STwoCreatMeanAIDramatic & \STwoCreatPenaltyDramaticEst & \STwoCreatPenaltyDramaticLo{} to \STwoCreatPenaltyDramaticHi & $p$~\STwoCreatPenaltyDramaticP\\
Averaged over levels\textsuperscript{a} & \STwoCreatMeanHuman & \STwoCreatMeanAI & \STwoCreatPenaltyEst & \STwoCreatPenaltyLo{} to \STwoCreatPenaltyHi & $p$~\STwoCreatPenaltyP\\
\bottomrule
\end{tabular}

\medskip
\begin{tabular}{@{}lcccccc@{}}
\multicolumn{7}{@{}l}{\textbf{b}\enspace Registered confirmatory tests of H1, with the ordinal check}\\
\toprule
 & & & & & \multicolumn{2}{c}{Ordinal check}\\
\cmidrule(l){6-7}
Step & Estimate & 95\% CI & $t$ ($df$) & After Holm & Log-odds & After Holm\\
\midrule
\multicolumn{7}{@{}l}{\textit{Transformativeness}}\\
Slight minus None & \STwoTransSlightMinusNoneEst & \STwoTransSlightMinusNoneLo{} to \STwoTransSlightMinusNoneHi & \STwoTransSlightMinusNoneT{} (\STwoTransSlightMinusNoneDf) & $p$~\STwoTransSlightMinusNonePHolm & \STwoTransSlightMinusNoneOrdEst & $p$~\STwoTransSlightMinusNoneOrdPHolm\\
Dramatic minus Slight & \STwoTransDramaticMinusSlightEst & \STwoTransDramaticMinusSlightLo{} to \STwoTransDramaticMinusSlightHi & \STwoTransDramaticMinusSlightT{} (\STwoTransDramaticMinusSlightDf) & $p$~\STwoTransDramaticMinusSlightPHolm & \STwoTransDramaticMinusSlightOrdEst & $p$~\STwoTransDramaticMinusSlightOrdPHolm\\
\addlinespace[3pt]
\multicolumn{7}{@{}l}{\textit{Essence change}}\\
Slight minus None & \STwoEssSlightMinusNoneEst & \STwoEssSlightMinusNoneLo{} to \STwoEssSlightMinusNoneHi & \STwoEssSlightMinusNoneT{} (\STwoEssSlightMinusNoneDf) & $p$~\STwoEssSlightMinusNonePHolm & \STwoEssSlightMinusNoneOrdEst & $p$~\STwoEssSlightMinusNoneOrdPHolm\\
Dramatic minus Slight & \STwoEssDramaticMinusSlightEst & \STwoEssDramaticMinusSlightLo{} to \STwoEssDramaticMinusSlightHi & \STwoEssDramaticMinusSlightT{} (\STwoEssDramaticMinusSlightDf) & $p$~\STwoEssDramaticMinusSlightPHolm & \STwoEssDramaticMinusSlightOrdEst & $p$~\STwoEssDramaticMinusSlightOrdPHolm\\
\addlinespace[3pt]
\multicolumn{7}{@{}l}{\textit{Creativity}}\\
Slight minus None & \STwoCreatSlightMinusNoneEst & \STwoCreatSlightMinusNoneLo{} to \STwoCreatSlightMinusNoneHi & \STwoCreatSlightMinusNoneT{} (\STwoCreatSlightMinusNoneDf) & $p$~\STwoCreatSlightMinusNonePHolm & \multicolumn{2}{c}{not completed\textsuperscript{b}}\\
Dramatic minus Slight & \STwoCreatDramaticMinusSlightEst & \STwoCreatDramaticMinusSlightLo{} to \STwoCreatDramaticMinusSlightHi & \STwoCreatDramaticMinusSlightT{} (\STwoCreatDramaticMinusSlightDf) & $p$~\STwoCreatDramaticMinusSlightPHolm & \multicolumn{2}{c}{}\\
\bottomrule
\end{tabular}

\smallskip
\begin{minipage}{\textwidth}
\footnotesize\raggedright
Predicted means and differences come from the registered model, averaged equally over image kind (originals with and without a human figure). A difference is computed before rounding, so it can differ in its last digit from the difference of the rounded means. $df$: degrees of freedom by Satterthwaite's approximation. The ordinal check is a cumulative-link mixed model, which treats the five answers as ordered categories. Its estimates are differences in the log-odds of a higher answer, and its Holm family is the nine ordinal $p$-values with the three equivalence $p$-values. \textsuperscript{a}The ordinary two-sided test of the averaged difference is description only; the registered H2 verdict rests on the equivalence test in Supplementary Table~\ref{tab:supp-labeleffects}. \textsuperscript{b}The ordinal model of creativity did not converge, so under the registered rule it has no $p$-value and no verdict.
\end{minipage}
\end{table}


\paragraph*{Registered text checks.}
Two checks on the creativity explanations were registered for every study, and Study~2 ran both (registered secondary; Section~\ref{sec:suppA-text} describes both). In the first, ridge regression, a regression that shrinks each word's weight toward zero, predicted each rating and each participant's label from the words of explanations written by participants left out of the fit. The words predicted the rating well: predicted and actual ratings correlated at \STwoTextRecRatingCreat{} (95\% confidence interval, CI, the range of values the data do not rule out at the 95\% level: \STwoTextRecRatingCreatLo{} to \STwoTextRecRatingCreatHi). In runs with the ratings shuffled, 95\% of correlations stayed below \STwoTextRecRatingCreatNull. The words did not show the label. The AUC, the chance that a randomly chosen AI-label explanation scores above a randomly chosen Human-label one, so that 0.5 is chance, was \STwoTextRecArmCreat{} (\STwoTextRecArmCreatLo{} to \STwoTextRecArmCreatHi), an interval that reaches chance; the rating alone gave \STwoTextRecArmFromRatingCreat. Adding the words to the rating raised the AUC by \STwoTextRecIncrCreat{} (\STwoTextRecIncrCreatLo{} to \STwoTextRecIncrCreatHi), which by the registered reading means that the words add to the rating. The second check predicted that AI-label participants would reuse their own words across explanations more than Human-label participants. It measured each person's word overlap among their own explanations, the share of words two explanations have in common, minus their overlap with other people's explanations of the same second work, with explanation length held constant. The difference, AI minus Human, was \STwoTextOverlapCreatDiff{} (\STwoTextOverlapCreatLo{} to \STwoTextOverlapCreatHi; Holm $p$~\STwoTextOverlapCreatPHolm), so the prediction is inconclusive, which is not evidence that the groups reuse words equally. Every explanation was written after its rating, so these checks show what the words track, not why anyone rated as they did.

\paragraph*{Exploratory text analyses.}
The registration names the remaining text analyses as exploratory, and none changes a verdict. Run on the transformativeness explanations, the first check told the label groups apart, though its interval cleared chance only barely; on the essence-change explanations it did not. Word tone and concreteness, scored from published norms, were compared between labels with a Benjamini--Hochberg correction, which limits the expected share of false positives among results called significant. A few small transformativeness terms survived it, among them slightly less pleasant words in AI-label explanations at the average rating; no creativity or essence-change term did. Topics found without labels, groups of words that tend to occur together, showed one transformativeness topic, led by the words `look', `identical' and `exact', with more weight in AI-label explanations after correction, and no creativity or essence-change topic that differed. A lasso, a regression that sets most word weights to exactly zero, kept a single word for telling the labels apart from creativity explanations: `ai'. Their estimates are in the Study~2 tables that the analysis code writes (main text, Code availability). The analyses pooled over Studies 2, 4 and 5 are in Supplementary Notes 6 and 7 and in Section~\ref{sec:suppB-crossstudy}.

\clearpage
